% ============================================================
% 第 6 章 结论与建议
% ============================================================

\section{结论与建议}
\label{sec:conclusion}

\subsection{主要结论}

本文围绕啤酒化合物—感官关系，完成了从数据对齐、分级评分、关键化合物识别到风格条件生成的三段式建模。主要结论可概括如下。

\textbf{问题一：}在缺少原始 0--6 等级标签的条件下，本文先用逐维均值/岭回归建立下界，再用缺失指示、分位数变换与风格先验改进多输出回归，最后用 PLS2、Ridge、ExtraTrees、RandomForest 与 LightGBM 的 OOF 等权融合与保序校准输出 0--6 分级。重复 $5\times3$ 折交叉验证显示，等权融合模型的宏平均 $R^2=0.227$、Spearman $=0.432$，保序校准后的 QWK 为 0.376；在 30\% MCAR 下 $R^2$ 相对下降 19.3\%，MAR 下下降 10.9\%，说明缺失指示与折内插补有效；保序校准把 0 级预测比例从 raw min-max 的 7.6\% 提到 11.3\%，明显缓解了分级塌缩，但仍低估参考分级的尾部（参考分级中 0 级占 23.3\%），这一限制在讨论中保留。

\textbf{问题二：}纯数据 Top-3 的平均 bootstrap 入选频率为 0.525、方法间 Jaccard 为 0.313；机制导向 Top-3 在 $\lambda=0.5$ 时把二乙酰、DMS、T2N 的直接化合物放回第 1/3/1 位。两者差异说明：在本数据中，缺陷直接化合物的边际信号弱于风格/物化代理，相关排序不能直接当作因果；共表达网络的 8 个社区与酒花萜类、酵母酯类、麦芽醛/呋喃/吡嗪对齐（$Q=0.586$），提供了独立于单变量相关的化学结构证据。

\textbf{问题三：}条件高斯 copula 与高斯最优传输在 log1p 空间建模；OT+约束在三风格上平均 MMD 为 0.038、覆盖度为 0.863，优于均值噪声（0.063/0.300）、全局协方差（0.066/0.233）与 copula（0.111/0.430）。Pils/Lager $\to$ Tripel 的 Tripel 分类命中率从 0.00 提升到 0.91；迁移后严格阈值下 DMS/Diacetyl/4VG/T2N 的超严格阈值率均为 0，style-p75 迁移在保留 Tripel 风格基线（DMS 39.8 \textmu g/L、Diacetyl 0.398 mg/L）的同时把 above-style-cap 率控制在 0；生成模型的 style cap 超标率四缺陷均为 0，严格阈值超标率按风格取舍另行报告。感官代理模型的嵌套 CV 平均 $R^2=0.220$，缺陷维度的预测较弱，这是本文明确报告的局限。

\subsection{对酿造实践的建议}

第一，把“化合物—感官”模型作为配方与工艺调整的辅助工具，而不是替代人工品评。模型可以快速筛出可能影响目标风味的化合物，但最终判断仍需结合感官小组与生产条件。

第二，对关键化合物建立分风格的监控窗口。问题二得到的 Top-3 可以作为优先监测项，但应结合问题一的维度权重与问题三的缺陷约束，避免只盯少数化合物而忽略整体风味平衡。

第三，缺陷控制应前置到原料与发酵过程。DMS、双乙酰、4VG、反式-2-壬烯醛与 stale hops 等指标与原料新鲜度、煮沸强度、发酵温度与储运条件密切相关；生成模型的阈值约束只能用于配方筛选，不能替代工艺控制。

第四，数据积累应优先补齐“缺失原因”和“原始 0--6 评分”。本文的校准层是对缺失标签的补救，不是对真实等级的替代；若能获得原始评分卡与检测限信息，问题一的绝对等级精度与问题三的感官评价都会显著改善。

\subsection{局限与展望}

本文的局限主要有四方面。其一，问题一缺少真实 0--6 标签，分级结果的绝对水平依赖校准假设；本文通过敏感性分析与有序约束降低风险，但不能消除。其二，问题二的关联分析基于观测数据，相关不等于因果；共表达图上的模块可能混合了共同调控、批次效应与测量误差。其三，问题三的生成质量评价缺少统一真值，本文的五维指标与案例只能提供相对比较，不能证明生成样本在真实酿造中可行。其四，外部数据只能整合到 22 个风格层级：Kaggle 数据与附件没有共同的酒名、酒厂或条码，无法逐酒连接；本文用它做风格级原型与外部效度校验，不能替代个体级标签\citep{kaggle2026topbeer}。

后续工作可以从三个方向展开：一是引入原始感官评分与检测限数据，把问题一的分级校准从“分位数锚定”升级为“有序标签监督”；二是用干预实验或代谢网络先验替代纯相关图，提升问题二因果解释的可信度；三是把问题三的生成模型与工艺参数连接，形成“工艺—化合物—感官”的闭环优化，而不是只在化合物层面做条件生成。
